% Part II -- Python Foundations
% Chapter 4: Booleans and Decisions

\h{Booleans and Decisions}

Programs often need to choose between actions. A room may pass or fail a size
check. A value may be inside or outside a permitted range. Python represents
these yes-or-no results with booleans and uses conditional statements to choose
which instructions to run.

\begin{NoteBox}
\textbf{Prerequisites.} You should understand variables, numbers, strings,
lists, dictionaries, membership, and basic input conversion.
\end{NoteBox}

\begin{tcolorbox}[title=Learning outcomes,colframe=orange,colback=white,arc=2mm]
After completing this chapter, you should be able to:
\begin{itemize}
    \item create boolean values with comparisons;
    \item combine conditions with \c{and}, \c{or}, and \c{not};
    \item write \c{if}, \c{elif}, and \c{else} statements;
    \item use membership tests in decisions; and
    \item trace which branch of a conditional statement will run.
\end{itemize}
\end{tcolorbox}

\hh{Boolean values}

Python has two boolean values: \c{True} and \c{False}. Their capital letters
matter.

\code{c04_boolean_values}{python}{
    is_accessible = True
    needs_review = False

    print(is_accessible)
    print(needs_review)
}{Store and display boolean values}

Boolean values usually come from questions expressed as comparisons.

\hh{Compare values}

Comparison operators produce \c{True} or \c{False}.

\begin{center}
\begin{tabular}{lll}
\textbf{Question} & \textbf{Operator} & \textbf{Example} \\
Equal to & \c{==} & \c{area_m2 == 30} \\
Not equal to & \c{!=} & \c{level != 1} \\
Greater than & \c{>} & \c{width_m > 3} \\
Less than & \c{<} & \c{width_m < 3} \\
Greater than or equal to & \c{>=} & \c{area_m2 >= 20} \\
Less than or equal to & \c{<=} & \c{area_m2 <= 40} \\
\end{tabular}
\end{center}

\code{c04_comparisons}{python}{
    area_m2 = 32.0

    print(area_m2 > 20.0)
    print(area_m2 == 32.0)
    print(area_m2 < 25.0)
}{Three comparisons}

Expected output:

\code{c04_comparisons_output}{text}{
    True
    True
    False
}{Comparison results}

Do not confuse \c{=} with \c{==}. A single equals sign assigns a value. Two
equals signs compare two values.

\begin{NoteBox}
\textbf{Stop and predict.} What are the results of \c{5 >= 5}, \c{5 > 5}, and
\c{"Studio" == "studio"}?
\end{NoteBox}

\hh{Check a range}

Python can express a range as one chained comparison.

\code{c04_chained_comparison}{python}{
    area_m2 = 32.0
    is_medium = 20.0 <= area_m2 <= 40.0

    print(is_medium)
}{Check whether a value is inside a range}

The expression is true when both boundaries are satisfied.

\hh{Combine conditions with and}

The \c{and} operator is true only when both conditions are true.

\code{c04_and_operator}{python}{
    width_m = 4.0
    length_m = 7.0

    dimensions_positive = width_m > 0 and length_m > 0
    print(dimensions_positive)
}{Require two conditions}

Use \c{and} when every requirement must be satisfied.

\hh{Combine alternatives with or}

The \c{or} operator is true when at least one condition is true.

\code{c04_or_operator}{python}{
    room_type = "Studio"

    is_work_space = room_type == "Studio" or room_type == "Office"
    print(is_work_space)
}{Accept either of two values}

Use \c{or} when one of several alternatives is enough.

\hh{Reverse a boolean with not}

The \c{not} operator reverses a boolean result.

\code{c04_not_operator}{python}{
    needs_review = False
    is_approved = not needs_review

    print(is_approved)
}{Reverse a boolean value}

Keep combined conditions short. A named boolean such as \c{is_approved} is
often easier to read than a long expression repeated several times.

\hh{Use membership conditions}

The \c{in} operator works with collections and returns a boolean.

\code{c04_membership_condition}{python}{
    allowed_types = ["Studio", "Office", "Classroom"]
    room_type = "Studio"

    is_allowed = room_type in allowed_types
    print(is_allowed)
}{Check whether a value is in a list}

Use \c{not in} to ask the opposite question.

\code{c04_not_in}{python}{
    allowed_types = ["Studio", "Office", "Classroom"]
    room_type = "Storage"

    print(room_type not in allowed_types)
}{Check that a value is absent}

\hh{Make a decision with if}

An \c{if} statement runs an indented block only when its condition is true.

\code{c04_first_if}{python}{
    area_m2 = 32.0

    if area_m2 >= 20.0:
        print("Area requirement passed")

    print("Check complete")
}{A basic if statement}

The colon ends the \c{if} line. The four spaces before the first \c{print()}
place that instruction inside the conditional block. The final \c{print()} is
not indented, so it runs regardless of the condition.

\begin{NoteBox}
\textbf{Stop and predict.} What will the program display if \c{area_m2} changes
to \c{18.0}?
\end{NoteBox}

\hh{Choose between two actions}

Use \c{if} and \c{else} when exactly one of two branches should run.

\code{c04_if_else}{python}{
    area_m2 = 18.0

    if area_m2 >= 20.0:
        print("Pass")
    else:
        print("Review")
}{Choose one of two branches}

Only one message appears. The \c{else} branch runs when the \c{if} condition is
false.

\hh{Choose among several actions}

Use \c{elif}, meaning ``else if,'' for additional conditions.

\code{c04_if_elif_else}{python}{
    area_m2 = 32.0

    if area_m2 < 20.0:
        category = "Small"
    elif area_m2 < 40.0:
        category = "Medium"
    else:
        category = "Large"

    print(category)
}{Classify a value with three branches}

Python checks the conditions from top to bottom and stops after the first true
branch. Because the first branch already handled values below 20, the second
branch only needs to check whether the area is below 40.

\hh{Read a decision program}

Trace values before running the code. Write the value of each variable, then
evaluate each condition as true or false.

\code{c04_trace_decision}{python}{
    width_m = 4.0
    length_m = 6.0
    area_m2 = width_m * length_m

    large_enough = area_m2 >= 20.0
    dimensions_positive = width_m > 0 and length_m > 0

    if large_enough and dimensions_positive:
        print("Accepted")
    else:
        print("Review the dimensions")
}{Trace named conditions before the decision}

Named conditions make the final decision easier to read and test.

\hh{Use basic truthiness carefully}

Python treats some values as false in a condition. These include \c{0}, an
empty string, and empty collections. Nonempty values usually count as true.

\code{c04_truthiness}{python}{
    room_name = "Studio"

    if room_name:
        print("A room name was provided")
    else:
        print("The room name is empty")
}{Use a string directly as a condition}

For beginners, explicit comparisons are often clearer. Use truthiness when the
meaning of empty versus nonempty is obvious.

\hh{Keep nesting shallow}

An \c{if} statement may appear inside another branch, but deep nesting becomes
difficult to follow.

\code{c04_nested_if}{python}{
    area_m2 = 32.0
    has_daylight = True

    if area_m2 >= 20.0:
        if has_daylight:
            print("Area and daylight checks passed")
        else:
            print("Daylight review required")
    else:
        print("Area review required")
}{A two-level nested decision}

When possible, combine related requirements with \c{and} or give conditions
clear names.

\hh{Common mistakes}

\begin{itemize}
    \item \textbf{Using \c{=} in a condition}: use \c{==} for comparison.
    \item \textbf{Forgetting the colon}: every \c{if}, \c{elif}, and \c{else}
          header ends with a colon.
    \item \textbf{Inconsistent indentation}: use four spaces for each block level.
    \item \textbf{Writing overlapping branches}: order conditions from specific
          or restrictive to general.
    \item \textbf{Comparing strings with inconsistent capitalization}: normalize
          or use the exact expected form.
\end{itemize}

\hh{Practice ladder}

\hhh{Level 0 -- Read and predict}

\code{c04_practice_level0}{python}{
    width_m = 3.0
    print(width_m > 2.5)
    print(width_m == 4.0)
}{Predict two boolean results}

\hhh{Level 1 -- Modify}

Run the following code, then try values below, inside, and above the range.

\code{c04_practice_level1}{python}{
    area_m2 = 25.0
    is_in_range = 20.0 <= area_m2 <= 40.0
    print(is_in_range)
}{Test several values against one range}

\hhh{Level 2 -- Complete}

Replace \c{...} with a condition that is true only when both dimensions are
positive.

\code{c04_practice_level2}{python}{
    width_m = 4.0
    length_m = 7.0

    dimensions_are_positive = ...

    if dimensions_are_positive:
        print("Pass")
    else:
        print("Review")
}{Complete a two-condition decision}

\hhh{Level 3 -- Apply}

Ask the user for an area. Display \c{Small} below 20, \c{Medium} from 20 up to
but not including 40, and \c{Large} at 40 or above.

\textbf{Hints}
\begin{itemize}
    \item \textbf{Hint 1}: Convert the input with \c{float()}.
    \item \textbf{Hint 2}: Use \c{if}, \c{elif}, and \c{else}.
    \item \textbf{Hint 3}: After the first condition fails, the value is already
          known to be at least 20.
\end{itemize}

\hhh{Level 4 -- Challenge}

Extend Level 3. If the area is zero or negative, display \c{Invalid area}
instead of a size category.

\hh{Practice solutions}

\textbf{Level 0:} The results are \c{True} and \c{False}.

\textbf{Level 2}

\code{c04_solution_level2}{python}{
    width_m = 4.0
    length_m = 7.0
    dimensions_are_positive = width_m > 0 and length_m > 0

    if dimensions_are_positive:
        print("Pass")
    else:
        print("Review")
}{One solution to Level 2}

\textbf{Level 3}

\code{c04_solution_level3}{python}{
    area_m2 = float(input("Area in m2: "))

    if area_m2 < 20.0:
        print("Small")
    elif area_m2 < 40.0:
        print("Medium")
    else:
        print("Large")
}{One solution to Level 3}

\textbf{Level 4}

\code{c04_solution_level4}{python}{
    area_m2 = float(input("Area in m2: "))

    if area_m2 <= 0:
        print("Invalid area")
    elif area_m2 < 20.0:
        print("Small")
    elif area_m2 < 40.0:
        print("Medium")
    else:
        print("Large")
}{One solution to Level 4}

\hh{Chapter checkpoint}

\begin{enumerate}
    \item What are Python's two boolean values?
    \item What is the difference between \c{=} and \c{==}?
    \item When is an \c{and} expression true?
    \item When is an \c{or} expression true?
    \item What does \c{not} do?
    \item What does \c{in} return?
    \item Which symbol ends an \c{if} header?
    \item How many spaces should indent one block level?
    \item When does an \c{else} branch run?
    \item Why does condition order matter in an \c{if--elif--else} statement?
\end{enumerate}

\hh{Checkpoint answers}

\begin{enumerate}
    \item \c{True} and \c{False}.
    \item \c{=} assigns; \c{==} compares.
    \item When both conditions are true.
    \item When at least one condition is true.
    \item It reverses a boolean result.
    \item A boolean indicating whether the value is present.
    \item A colon.
    \item Four spaces.
    \item When the preceding conditions are false.
    \item Python runs only the first matching branch.
\end{enumerate}

\begin{tcolorbox}[title=Chapter summary,colframe=orange,colback=white,arc=2mm]
Comparisons produce boolean values. Logical operators combine or reverse those
results. Conditional statements use the results to choose which instructions
run. Chapter 5 uses loops to repeat instructions for several values.
\end{tcolorbox}

